
\section{Conclusion}
\label{sec:conclusion}
In this paper, we introduce Qwen3-Omni-30B-A3B, Qwen3-Omni-30B-A3B-Thinking, Qwen3-Omni-Flash-Instruct, and Qwen3-Omni-Flash-Thinking models. Qwen3-Omni-30B-A3B matches or surpasses the latest same-size unimodal Qwen models on text and vision benchmarks. Notably, on audio processing and dialogue benchmarks, it attains state-of-the-art performance among open-source systems on 32 benchmarks and is comparable to, or better than, the strong proprietary counterpart Gemini-2.5-Pro. The Qwen3-Omni-30B-A3B Thinking variant achieves further gains on complex tasks spanning text, vision, and audio-visual reasoning. Beyond accuracy, the model supports 119 text languages, 19 languages for speech recognition and 10 languages for speech synthesis, and enables audio understanding and interactive sessions up to 40 minutes. Thanks to its streaming architecture and multi-codebook design, Qwen3-Omni at the 30B-A3B scale still delivers an end-to-end first-packet latency of 234 ms.

Research fields often cycle between specialization and integration. In this context, we believe Qwen3-Omni represents a milestone: to our knowledge, it provides the first evidence that fully integrated, end-to-end multimodal training can be achieved without degrading core language capability and other modalities. We are eager to share these findings with the community and hope they will stimulate further research.

For practical usage, Qwen3-Omni-30B-A3B offers strong text and vision capabilities, robust and reliable ASR, interactive speech support in over 20 languages, very low first-packet latency for interactive use, and stable, naturalistic speech synthesis. Crucially, it exhibits advantages over cascaded pipelines, including stronger cross-modal reasoning, lower end-to-end latency, and lower system complexity and cost. In future work, we will further advance the model along multiple axes, including multi-speaker ASR, video OCR, audiovisual proactive learning, and enhanced support for agent-based workflows and function calling.
